\subsection{Pfaffian of a skew-symmetric matrix.}

Let A be a commutative field of characteristic 0, and \(R = (\alpha_{ij})\) an alternating matrix of even order \(2m\) on A. Let E denote the vector space \(\mathbf{A}^{2m}\) , \((e_i)\) ( \(i = 1, \ldots, 2m\) ) its canonical basis, and \(e\) the element \(e_1 \wedge e_2 \wedge \ldots \wedge e_{2m}\) of \(\bigwedge^{2m}\) E. The \(m\) -th exterior power of the bivector \(u = \sum_{i < j} \alpha_{ij} e_i \wedge e_j \in \bigwedge^2 E\) is of the form \(\alpha.e\) , where \(\alpha\) is an element of A that we will compute. The element \(\bigwedge^m u\) is a sum of terms of the form

\[
\alpha_ {h _ {1} k _ {1}} \alpha_ {h _ {2} k _ {2}} \dots \alpha_ {h _ {m} k _ {m}} e _ {h _ {1}} \wedge e _ {k _ {1}} \wedge e _ {h _ {2}} \wedge e _ {k _ {2}} \wedge \dots \wedge e _ {h _ {m}} \wedge e _ {k _ {m}}\tag{3}
\]

with \(h_j < k_j\) for \(j = 1, \ldots, m\) . Such a term is zero if it contains two equal \(e_j\) , that is, if the set \(\{h_1, k_1, \ldots, h_m, k_m\}\) is not exactly \(\{1, 2, \ldots, 2m\}\) . Moreover, if, in (3), one simultaneously interchanges \(e_{h_r}\) and \(e_{h_{r+1}}\) on the one hand, \(e_{k_r}\) and \(e_{k_{r+1}}\) on the other hand, the product does not change; hence it does not change under any permutation performed on the pairs \((h_1, k_1), \ldots, (h_m, k_m)\) . Consider then the sets (and not the sequences) \(S = \{(h_1, k_1), \ldots, (h_m, k_m)\}\) of pairs \((h_j, k_j)\) such that \(1 \leqslant h_j < k_j \leqslant 2m\) for \(j = 1, 2, \ldots, m\) ; let \(\mathscr{F}\) be the set of these pairs. For \(S \in \mathscr{F}\) , set

\[
1 ^ {0}) \varepsilon (S) = 0 \text {   si   } \left\{h _ {1}, k _ {1}, \dots , h _ {m}, k _ {m} \right\} \neq \left\{1, 2, \dots , 2 m \right\};
\]

2°) in the contrary case, \(\varepsilon(S)=1\) or \(\varepsilon(S)=-1\) depending on whether the permutation that maps \(h_{j}\) to 2j-1 and \(k_{j}\) to 2j ( \(j=1,\ldots,m\) ) is even or odd.

The preceding remarks then prove that \(\overset{m}{\wedge} u\) is equal to

\[
m! \sum_ {S \in \mathfrak {S}} \varepsilon (S) (\prod_ {(h, k) \in S} \alpha_ {h k}) e.\tag{4}
\]

We now introduce \(m(2m-1)\) indeterminates \(X_{hk}\) indexed by the pairs \((h,k)\) such that \(1\leqslant h<k\leqslant2m\) and let us call P the polynomial over \(\mathbf{Z}\) with respect to the \(X_{hk}\) defined by

\[
\mathrm{P} ((X _ {h k})) = \sum_ {\mathbf {S} \in \mathfrak {S}} \varepsilon (\mathrm{S}) (\prod_ {(h, k) \in \mathrm{S}} X _ {h k}).\tag{5}
\]

We therefore have

\[
\stackrel {m} {\wedge} u = m! \mathrm{P} ((\alpha_ {h k})). e.\tag{6}
\]

DEFINITION 1. --- Given an alternating matrix \(R = (\alpha_{ij}) (i, j = 1, \ldots, 2m)\) of even order \(2m\) on an arbitrary commutative ring A, one calls the Pfaffian of \(R\) , and denotes it \(\mathrm{Pf}(R)\) , the element \(\mathrm{P}((\alpha_{hk}))\) of A, where \(1 \leqslant h < k \leqslant 2m\) .

Example. --- Suppose:

\[
\alpha_ {1 2} = - \alpha_ {2 1} = \beta_ {1}, \alpha_ {3 4} = - \alpha_ {4 3} = \beta_ {2}, \dots , \alpha_ {2 m - 1, 2 m} = - \alpha_ {2 m, 2 m - 1} = \beta_ {m},
\]

all the others \(\alpha_{ij}\) being zero (cf. th. 1). Then the Pfaffian of \(R = (\alpha_{ij})\) is \(\beta_1\beta_2\ldots\beta_m\) .

PROPOSITION 1. --- Let R be an alternating matrix of even order 2m over a commutative ring A, and let P be a square matrix of order 2m over A. Then

\[
\operatorname{Pf} (^ {t} P. R. P) = (\det P) \operatorname{Pf} (R).\tag{7}
\]

Indeed, suppose first that A is a field of characteristic 0 and set \(R = (\alpha_{ij})\) , \(P = (\beta_{st})\) . Associate with \(R\) the bivector

\[
u = \sum_ {i <   j} \alpha_ {i j} e _ {i} \wedge e _ {j} = \frac {1}{2} \sum_ {1 \leqslant i, j \leqslant 2 m} \alpha_ {i j} e _ {i} \wedge e _ {j}
\]

of \(\wedge\) A \(^{2m}\) , where (e \(_i\) ) denotes the canonical basis of A \(^{2m}\) ; consider \(^t P\) as the matrix, with respect to the basis (e \(_i\) ), of an endomorphism \(f\) of A \(^{2m}\) . Then the bivector ( \(\overset{2}{\wedge} f\) )(u) is associated with the matrix \(^t P.R.P\) since it is equal to \(\frac{1}{2}\sum_{i,j,s,t}\beta_{is}\alpha_{ij}\beta_{jt}e_s\wedge e_t\) . Since the extension \(\wedge f\) of \(f\) to the exterior algebra \(\wedge\) A \(^{2m}\) is an endomorphism of this algebra (chap. III, § 5, no. 9), we have \(\overset{m}{\wedge}((\overset{2}{\wedge}f)(u)) = (\overset{2m}{\wedge}f)(\overset{m}{\wedge}u)\) ; since \(\overset{2m}{\wedge} f\) is the homothety with ratio det \(f\) , it therefore follows from (6) and from def. 1 that we have \(m!\) Pf( \(^t P.R.P\) ) = \(m!\) (det \(P\) )Pf( \(R\) ), whence (7) in the case under consideration. The general case follows by noting that both sides of (7) are polynomials with integer coefficients in the entries of the matrices \(R\) and \(P\) (chap. IV, § 2, no. 5, Scholium).

PROPOSITION 2. --- For every alternating matrix R of even order 2m over a commutative ring A, we have

\[
\det R = (\operatorname{Pf} (R)) ^ {2}.\tag{8}
\]

Indeed, since both sides of (8) are polynomials with integer coefficients in the elements of \(R\) , the principle of extension of algebraic identities (chap. IV, § 2, no. 5, Scholium) shows that it suffices to carry out the proof in the case where A is a field of characteristic 0 and where det \(R \neq 0\) . If \(P\) is an invertible square matrix of order \(2m\) on A, we have det( \(^{t}P.R.P\) ) = ( \(\det P\) ) \(^{2}\) ( \(\det R\) ) and \(\mathrm{Pf}(^{t}P.R.P)\) = ( \(\det P\) ) \(\mathrm{Pf}(R)\) (prop. 1), so that it suffices to prove (8) for \(^{t}P.R.P\) instead of \(R\) . By cor. 1 of th. 1, one can, by a suitable choice of \(P\) , suppose that the matrix \(^{t}P.R.P\) is of the form ( \(\alpha_{ij}\) ) where

\[
\alpha_ {1 2} = - \alpha_ {2 1} = 1, \dots , \alpha_ {2 m - 1, 2 m} = - \alpha_ {2 m, 2 m - 1} = 1,
\]

all the others \(\alpha_{ij}\) being zero (cf. Example). Now the determinant of this matrix equals 1, and its Pfaffian likewise; this completes the proof.

